\chapter{控制器实现代码附录}\label{chap:control-code}
\subsection{控制器部分}
\begin{minted}[breaklines,linenos,fontsize=\small]{C++}

namespace controllers
{

TandemLeg::TandemLeg(
  sp::Mahony & imu, float wheel_radius, float wheel_distance, float l1, float l2, float l3,
  float l4, float lf_ecd, float lr_ecd, float rf_ecd, float rr_ecd, bool reverse_lf,
  bool reverse_lr, bool reverse_rf, bool reverse_rr, bool reverse_left_wheel,
  bool reverse_right_wheel)
: imu_(imu),
  wheel_radius_(wheel_radius),
  wheel_distance_(wheel_distance),
  l1_(l1),
  l2_(l2),
  l3_(l3),
  l4_(l4),
  lf_ecd_(lf_ecd),
  lr_ecd_(lr_ecd),
  rf_ecd_(rf_ecd),
  rr_ecd_(rr_ecd),
  reverse_lf_(reverse_lf),
  reverse_lr_(reverse_lr),
  reverse_rf_(reverse_rf),
  reverse_rr_(reverse_rr),
  reverse_left_wheel_(reverse_left_wheel),
  reverse_right_wheel_(reverse_right_wheel)
{
  this->bar.l1 = l1_;
  this->bar.l2 = l2_;
  this->bar.l3 = l3_;
  this->bar.l4 = l4_;
  this->bar.l5 = 0.0f;
}

void TandemLeg::disable() { mode_ = ControlMode::DISABLE; }

//普通力控模式
void TandemLeg::cmd_force(ChassisState & cmd)
{
  mode_ = ControlMode::FORCE;
  cmd_ = cmd;
  //关闭s项和rotate项
  cmd_.s = fdb_.s;
  cmd_.rotate = fdb_.rotate;
}

//翻倒起身速控
void TandemLeg::cmd_velocity() { mode_ = ControlMode::VELOCITY; }

//防溜车“位控”
void TandemLeg::cmd_position(ChassisState & cmd)
{
  mode_ = ControlMode::POSITION;
  cmd_ = cmd;
}

void TandemLeg::init()
{
  left_wheel_ecd_ = chassis::motor_left_wheel.angle;
  right_wheel_ecd_ = chassis::motor_right_wheel.angle;
}

void TandemLeg::control()
{
  if (mode_ == ControlMode::DISABLE) {
    chassis::motor_lr.cmd(0.0f);
    chassis::motor_lf.cmd(0.0f);
    chassis::motor_rr.cmd(0.0f);
    chassis::motor_rf.cmd(0.0f);
    chassis::motor_left_wheel.cmd(0.0f);
    chassis::motor_right_wheel.cmd(0.0f);
    return;
  }
  if (mode_ == ControlMode::VELOCITY) {
  }
  if (mode_ == ControlMode::FORCE) {
#ifdef K_FIT_ENABLED
    lqr.k_fit(tandem_leg.leg_l.l0, tandem_leg.leg_r.l0);
#endif
    tandem_leg.input = lqr.calc(cmd_, fdb_);
    tandem_leg.F0_calc();
    tandem_leg.vmc_calc();
  }
  if (mode_ == ControlMode::POSITION) {
#ifdef K_FIT_ENABLED
    lqr.k_fit(tandem_leg.leg_l.l0, tandem_leg.leg_r.l0);
#endif
    tandem_leg.input = lqr.calc(cmd_, fdb_, true);
    tandem_leg.F0_calc();
    tandem_leg.vmc_calc();
  }
  chassis::motor_lf.cmd(sp::limit_max(tandem_leg.input.t_lf, TandemLeg::MAX_HIP_TORQUE));
  chassis::motor_lr.cmd(sp::limit_max(tandem_leg.input.t_lr, TandemLeg::MAX_HIP_TORQUE));
  chassis::motor_rf.cmd(sp::limit_max(tandem_leg.input.t_rf, TandemLeg::MAX_HIP_TORQUE));
  chassis::motor_rr.cmd(sp::limit_max(tandem_leg.input.t_rr, TandemLeg::MAX_HIP_TORQUE));
  chassis::motor_left_wheel.cmd(sp::limit_max(tandem_leg.input.t_wl, TandemLeg::MAX_WHEEL_TORQUE));
  chassis::motor_right_wheel.cmd(sp::limit_max(tandem_leg.input.t_wr, TandemLeg::MAX_WHEEL_TORQUE));
  // chassis::motor_lr.cmd(0.0f);
  // chassis::motor_lf.cmd(0.0f);
  // chassis::motor_rr.cmd(0.0f);
  // chassis::motor_rf.cmd(0.0f);
  // chassis::motor_left_wheel.cmd(0.0f);
  // chassis::motor_right_wheel.cmd(0.0f);
}

void TandemLeg::update()
{
  chassis::filter_chassis();
  update_wheel(
    chassis::left_wheel_speed_filter.out, chassis::right_wheel_speed_filter.out,
    chassis::left_wheel_angle_filter.out, chassis::right_wheel_angle_filter.out, left_wheel_ecd_,
    right_wheel_ecd_);
  update_joint(
    chassis::motor_lf.angle, chassis::motor_lr.angle, chassis::motor_rf.angle,
    chassis::motor_rr.angle, chassis::lf_speed_filter.out, chassis::lr_speed_filter.out,
    chassis::rf_speed_filter.out, chassis::rr_speed_filter.out);
  update_leg(leg_l);
  update_leg(leg_r);
  update_chassis_state();
  //KF观测速度
  kf_update_z();
  ekf::kf.predict(ekf::F, ekf::Q);
  // ekf::kf.update(ekf::z, ekf::H, ekf::R);

  slip_detect();
  support_F_calc();
}

void TandemLeg::update_wheel(
  float left_wheel_speed, float right_wheel_speed, float left_wheel_angle, float right_wheel_angle,
  float left_wheel_ecd, float right_wheel_ecd)
{
  state.l_dtheta = (reverse_left_wheel_) ? -left_wheel_speed : left_wheel_speed;
  state.r_dtheta = (reverse_right_wheel_) ? -right_wheel_speed : right_wheel_speed;
  state.l_theta = (reverse_left_wheel_ ? -1.0f : 1.0f) * (left_wheel_angle - left_wheel_ecd);
  state.r_theta = (reverse_right_wheel_ ? -1.0f : 1.0f) * (right_wheel_angle - right_wheel_ecd);
}

void TandemLeg::update_joint(
  float lf_angle, float lr_angle, float rf_angle, float rr_angle, float lf_speed, float lr_speed,
  float rf_speed, float rr_speed)
{
  state.l_phi1 = sp::limit_angle((reverse_lf_ ? -1.0f : 1.0f) * (lf_angle - this->lf_ecd_));
  state.l_phi4 = sp::limit_angle((reverse_lr_ ? -1.0f : 1.0f) * (lr_angle - this->lr_ecd_));
  state.r_phi1 = sp::limit_angle((reverse_rf_ ? -1.0f : 1.0f) * (rf_angle - this->rf_ecd_));
  state.r_phi4 = sp::limit_angle((reverse_rr_ ? -1.0f : 1.0f) * (rr_angle - this->rr_ecd_));

  state.l_dphi1 = (reverse_lf_ ? -1.0f : 1.0f) * lf_speed;
  state.l_dphi4 = (reverse_lr_ ? -1.0f : 1.0f) * lr_speed;
  state.r_dphi1 = (reverse_rf_ ? -1.0f : 1.0f) * rf_speed;
  state.r_dphi4 = (reverse_rr_ ? -1.0f : 1.0f) * rr_speed;

  leg_l.phi1 = state.l_phi1;
  leg_l.phi4 = state.l_phi4;
  leg_r.phi1 = state.r_phi1;
  leg_r.phi4 = state.r_phi4;
}

void TandemLeg::update_leg(Leg & leg)
{
  const float L5 = 0.0f;

  float YD = l4_ * std::sin(leg.phi4);
  float YB = l1_ * std::sin(leg.phi1);
  float XD = L5 + l4_ * std::cos(leg.phi4);
  float XB = l1_ * std::cos(leg.phi1);

  float lBD = std::sqrt((XD - XB) * (XD - XB) + (YD - YB) * (YD - YB));
  float A0 = 2 * l2_ * (XD - XB);
  float B0 = 2 * l2_ * (YD - YB);
  float C0 = l2_ * l2_ + lBD * lBD - l3_ * l3_;

  leg.phi2 = 2 * std::atan2((B0 + std::sqrt(A0 * A0 + B0 * B0 - C0 * C0)), (A0 + C0));
  leg.phi3 = std::atan2(YB - YD + l2_ * std::sin(leg.phi2), XB - XD + l2_ * std::cos(leg.phi2));

  float XC = l1_ * std::cos(leg.phi1) + l2_ * std::cos(leg.phi2);
  float YC = l1_ * std::sin(leg.phi1) + l2_ * std::sin(leg.phi2);

  float L0 = std::sqrt((XC - L5 / 2) * (XC - L5 / 2) + YC * YC);
  float phi0 = std::atan2(YC, (XC - L5 / 2));

  leg.xc = XC;
  leg.yc = YC;
  leg.l0 = L0;
  leg.phi0 = phi0;

  leg.phi0_real = phi0 + imu_.pitch;

  leg.dphi2 = 1 / (-l2_ * sin(leg.phi2 - leg.phi3)) *
              (l1_ * sin(leg.phi1 - leg.phi3) * state.l_dphi1 -
               l4_ * sin(leg.phi4 - leg.phi3) * state.l_dphi4);
  if (&leg == &leg_l) {
    fdb_.theta_ll = leg.phi0_real - sp::Pi / 2.0f;
    fdb_.l_l0 = leg.l0;
  }
  else if (&leg == &leg_r) {
    fdb_.theta_lr = leg.phi0_real - sp::Pi / 2.0f;
    fdb_.r_l0 = leg.l0;
  }
  get_fdb_dtheta();
  get_fdb_dl0();
}

void TandemLeg::update_chassis_state()
{
  if (mode_ == ControlMode::POSITION) {
    this->fdb_.s += this->fdb_.ds * chassis::CHASSIS_CONTROL_T;
    this->fdb_.rotate += this->fdb_.drotate * chassis::CHASSIS_CONTROL_T;
  }
  else {
    this->fdb_.s = 0.0f;
    this->fdb_.rotate = 0.0f;
  }
  this->fdb_.ds = this->wheel_radius_ / 2.0f * (this->state.l_dtheta + this->state.r_dtheta);
  this->fdb_.drotate =
    this->wheel_radius_ / 2.0f / this->wheel_distance_ *
      (this->state.r_dtheta - this->state.l_dtheta) -
    leg_l.l0 / 2.0f / this->wheel_distance_ * cos(fdb_.theta_ll) * fdb_.dtheta_ll +
    leg_r.l0 / 2.0f / this->wheel_distance_ * cos(fdb_.theta_lr) * fdb_.dtheta_lr;
  this->fdb_.phi = imu_.pitch;
  this->fdb_.dphi = imu_.vpitch;
}

void TandemLeg::vmc_calc()
{
  //左腿
  float l_j11 = (bar.l1 * sin(leg_l.phi0 - leg_l.phi3) * sin(leg_l.phi1 - leg_l.phi2)) /
                sin(leg_l.phi3 - leg_l.phi2);
  float l_j12 = (bar.l1 * cos(leg_l.phi0 - leg_l.phi3) * sin(leg_l.phi1 - leg_l.phi2)) /
                (leg_l.l0 * sin(leg_l.phi3 - leg_l.phi2));
  float l_j21 = (bar.l4 * sin(leg_l.phi0 - leg_l.phi2) * sin(leg_l.phi3 - leg_l.phi4)) /
                sin(leg_l.phi3 - leg_l.phi2);
  float l_j22 = (bar.l4 * cos(leg_l.phi0 - leg_l.phi2) * sin(leg_l.phi3 - leg_l.phi4)) /
                (leg_l.l0 * sin(leg_l.phi3 - leg_l.phi2));

  input.t_lf = (reverse_lf_ ? -1.0f : 1.0f) * (l_j11 * input.f0_l + l_j12 * input.tp_l);
  input.t_lr = (reverse_lr_ ? -1.0f : 1.0f) * (l_j21 * input.f0_l + l_j22 * input.tp_l);
  input.t_wl = (reverse_left_wheel_ ? -1.0f : 1.0f) * input.t_wl;  //左轮

  //右腿
  float r_j11 = (bar.l1 * sin(leg_r.phi0 - leg_r.phi3) * sin(leg_r.phi1 - leg_r.phi2)) /
                sin(leg_r.phi3 - leg_r.phi2);
  float r_j12 = (bar.l1 * cos(leg_r.phi0 - leg_r.phi3) * sin(leg_r.phi1 - leg_r.phi2)) /
                (leg_r.l0 * sin(leg_r.phi3 - leg_r.phi2));
  float r_j21 = (bar.l4 * sin(leg_r.phi0 - leg_r.phi2) * sin(leg_r.phi3 - leg_r.phi4)) /
                sin(leg_r.phi3 - leg_r.phi2);
  float r_j22 = (bar.l4 * cos(leg_r.phi0 - leg_r.phi2) * sin(leg_r.phi3 - leg_r.phi4)) /
                (leg_r.l0 * sin(leg_r.phi3 - leg_r.phi2));

  input.t_rf = (reverse_rf_ ? -1.0f : 1.0f) * (r_j11 * input.f0_r + r_j12 * input.tp_r);
  input.t_rr = (reverse_rr_ ? -1.0f : 1.0f) * (r_j21 * input.f0_r + r_j22 * input.tp_r);
  input.t_wr = (reverse_right_wheel_ ? -1.0f : 1.0f) * input.t_wr;  //右轮
}

void TandemLeg::F0_calc()
{
  //支持力=重力补偿+腿长补偿+roll补偿+侧向惯性力矩补偿
  float chassis_force = chassis::CHASSIS_WEIGHT * chassis::GRAVITY_ACC;
  //重力补偿
  float F0_L = chassis_force * cosf(fdb_.theta_ll) / 2.0f;
  float F0_R = chassis_force * cosf(fdb_.theta_lr) / 2.0f;
  //腿长补偿
  l_leg_length_pid.calc(cmd_.l_l0, leg_l.l0);
  r_leg_length_pid.calc(cmd_.l_l0, leg_r.l0);
  //roll补偿
  leg_roll_pid.calc(0.0f, imu_.roll);
  //侧向惯性力矩补偿
  float lateral_inertia_torque =
    chassis::CHASSIS_WEIGHT / 2 * 1 / (2 * chassis::WHEEL_DISTANCE) * fdb_.dphi * fdb_.s;
  input.f0_l = F0_L + l_leg_length_pid.out + leg_roll_pid.out - lateral_inertia_torque;
  input.f0_r = F0_R + r_leg_length_pid.out - leg_roll_pid.out + lateral_inertia_torque;
}

void TandemLeg::get_fdb_dtheta()
{
  get_fdb_dtheta_leg(leg_l);
  get_fdb_dtheta_leg(leg_r);
}

void TandemLeg::get_fdb_dtheta_leg(Leg & leg)
{
  float J11 = bar.l1 * sin(leg.phi1 - leg.phi2) * sin(leg.phi3) / sin(leg.phi2 - leg.phi3);
  float J12 = bar.l4 * sin(leg.phi3 - leg.phi4) * sin(leg.phi2) / sin(leg.phi2 - leg.phi3);
  float J21 = -bar.l1 * sin(leg.phi1 - leg.phi2) * sin(leg.phi3) / sin(leg.phi2 - leg.phi3);
  float J22 = -bar.l4 * sin(leg.phi3 - leg.phi4) * sin(leg.phi2) / sin(leg.phi2 - leg.phi3);

  if (&leg == &leg_l) {
    float dxc_l = J11 * state.l_dphi1 + J12 * state.l_dphi4;
    float dyc_l = J21 * state.l_dphi1 + J22 * state.l_dphi4;
    fdb_.dtheta_ll =
      -(leg.yc * dxc_l - leg.xc * dyc_l) / (leg.xc * leg.xc + leg.yc * leg.yc) + imu.vpitch;
  }
  if (&leg == &leg_r) {
    float dxc_l = J11 * state.r_dphi1 + J12 * state.r_dphi4;
    float dyc_l = J21 * state.r_dphi1 + J22 * state.r_dphi4;
    fdb_.dtheta_lr =
      -(leg.yc * dxc_l - leg.xc * dyc_l) / (leg.xc * leg.xc + leg.yc * leg.yc) + imu.vpitch;
  }
}

void TandemLeg::get_fdb_dl0()
{
  get_fdb_dl0_leg(leg_l);
  get_fdb_dl0_leg(leg_r);
}

void TandemLeg::get_fdb_dl0_leg(Leg & leg)
{
  if (&leg == &leg_l) {
    float dXC = -bar.l1 * sin(leg.phi1) * state.l_dphi1 - bar.l2 * sin(leg.phi2) * leg_l.dphi2;
    float dYC = bar.l1 * cos(leg.phi1) * state.l_dphi1 + bar.l2 * cos(leg.phi2) * leg_l.dphi2;
    fdb_.dl_l0 = (leg.xc * dXC + leg.yc * dYC) / leg.l0;
  }
  if (&leg == &leg_r) {
    float dXC = -bar.l1 * sin(leg.phi1) * state.r_dphi1 - bar.l2 * sin(leg.phi2) * leg_r.dphi2;
    float dYC = bar.l1 * cos(leg.phi1) * state.r_dphi1 + bar.l2 * cos(leg.phi2) * leg_r.dphi2;
    fdb_.dr_l0 = (leg.xc * dXC + leg.yc * dYC) / leg.l0;
  }
}

//KF更新
//需要在imu已经跑起来后调用
void TandemLeg::kf_update_z()
{
  static float last_gyro_z = imu_.vyaw;
  float acc_z = (imu_.vyaw - last_gyro_z) / CONTROL_T;

  ekf::z[0][0] = fdb_.ds;
  ekf::z[1][0] = bmi088.acc[0] * cosf(imu.pitch) - chassis::GRAVITY_ACC * sinf(imu.pitch);
  ekf::z[2][0] = fdb_.drotate;
  ekf::z[3][0] = acc_z;
  if (remote.sw_l == sp::DBusSwitchMode::DOWN) {
    ekf::kf_initialized = false;
  }
  if (!ekf::kf_initialized) {
    ekf::kf_initialized = true;
    ekf::kf.x[0][0] = ekf::z[0][0];
    ekf::kf.x[1][0] = ekf::z[1][0];
    ekf::kf.x[2][0] = ekf::z[2][0];
    ekf::kf.x[3][0] = ekf::z[3][0];
    return;
  }
}

void TandemLeg::slip_detect()
{
  // float innovation_v = chassis.fdb.ds - (z[1][0] * T_CONTROL +);  // 残差

  // 核心逻辑：如果残差过大，说明编码器速度突变（打滑），
  // 此时我们不信任编码器，极大幅度增加测量噪声 R
  if (std::abs(ekf::kf.data["ds"]) > SLIP_THRESHOLD) {
    // 检测到打滑
    // 策略：瞬间增大测量噪声，让滤波器主要保留预测值（即保留 IMU 积分结果）
    ekf::R = ekf::R_slip_data;
    slip_flag_ = true;
    // Q = q_slip_diag_vec;
    ekf::Q = ekf::Q_data_slip;
    ekf::Q[0][0] = 0.1f;
    ekf::Q[1][1] = 10.0f;
    ekf::Q[2][2] = 0.5f;
    ekf::Q[3][3] = 0.5f;
  }
  // else {
  //   R = R_data;
  //   slip_flag_ = false;
  // }
  //在0附近时来回切？
  // if(std::abs(innovation_v) <= SLIP_OVER_THRESHOLD) {
  //     // 正常情况，恢复正常测量噪声
  //     R = R_data;
  // }
}

void TandemLeg::support_F_calc()
{
  support_F_calc_leg(leg_l);
  support_F_calc_leg(leg_r);
}

void TandemLeg::support_F_calc_leg(Leg & leg)
{
  float zM = ins_accel_corrected[2];
  if (&leg == &leg_l) {
    float P = input.f0_l * cosf(state.l_theta) + input.tp_l * sinf(state.l_theta) / leg.l0;
    state.l_F0 = P + chassis::CHASSIS_WEIGHT * chassis::GRAVITY_ACC / 2.0f +
                 chassis::CHASSIS_WEIGHT * zM / 2.0f;
  }
  if (&leg == &leg_r) {
    float P = input.f0_r * cosf(state.r_theta) + input.tp_r * sinf(state.r_theta) / leg.l0;
    state.r_F0 = P + chassis::CHASSIS_WEIGHT * chassis::GRAVITY_ACC / 2.0f +
                 chassis::CHASSIS_WEIGHT * zM / 2.0f;
  }
}
}  // namespace controllers


\end{minted}

\subsection{交互部分}
\begin{minted}[breaklines,linenos,fontsize=\small]{python}
import math

import matplotlib.pyplot as plt
import numpy as np

show_animation = True  # 动画


class Config(object):
    """
    用来仿真的参数，
    """

    def __init__(self):
        # robot parameter
        self.max_speed = 0.5  # [m/s]  # 最大速度
        self.min_speed = -0.5  # [m/s]  # 最小速度，设置为可以倒车
        # self.min_speed = 0  # [m/s]  # 最小速度，设置为不倒车
        self.max_yawrate = 2  # [rad/s]  # 最大角速度
        self.max_accel = 0.2  # [m/ss]  # 最大加速度
        self.max_dyawrate = 2  # [rad/ss]  # 最大角加速度
        self.v_reso = 0.01  # [m/s]，速度分辨率
        self.yawrate_reso = 0.1  # [rad/s]，角速度分辨率
        self.dt = 0.1  # [s]  # 采样周期
        self.predict_time = 3.0  # [s]  # 向前预估三秒
        self.to_goal_cost_gain = 1.0  # 目标代价增益
        self.speed_cost_gain = 1.0  # 速度代价增益
        self.robot_radius = 0.09  # [m]  # 机器人半径

def motion(x, u, dt):
    """
    :param x: 位置参数，在此叫做位置空间
    :param u: 速度和加速度，在此叫做速度空间
    :param dt: 采样时间
    :return:
    """
    x[0] += u[0] * math.cos(x[2]) * dt  # x方向位移
    x[1] += u[0] * math.sin(x[2]) * dt  # y方向位移
    x[2] += u[1] * dt  # 航向角
    x[3] = u[0]  # 速度v
    x[4] = u[1]  # 角速度w
    return x


def calc_dynamic_window(x, config):
    """
    位置空间集合
    :param x:当前位置空间,符号参考硕士论文
    :param config:
    :return:目前是两个速度的交集，还差一个
    """

    # 车辆能够达到的最大最小速度
    vs = [config.min_speed, config.max_speed, -config.max_yawrate, config.max_yawrate]

    # 一个采样周期能够变化的最大最小速度
    vd = [
        x[3] - config.max_accel * config.dt,
        x[3] + config.max_accel * config.dt,
        x[4] - config.max_dyawrate * config.dt,
        x[4] + config.max_dyawrate * config.dt,
    ]
    #  print(Vs, Vd)

    # 求出两个速度集合的交集
    vr = [max(vs[0], vd[0]), min(vs[1], vd[1]), max(vs[2], vd[2]), min(vs[3], vd[3])]

    return vr


# 产生运动轨迹
def calc_trajectory(x_init, v, w, config):
    """
    预测3秒内的轨迹
    :param x_init:位置空间
    :param v:速度
    :param w:角速度
    :param config:
    :return: 每一次采样更新的轨迹，位置空间垂直堆叠
    """
    x = np.array(x_init)
    trajectory = np.array(x)
    time = 0
    while time <= config.predict_time:
        x = motion(x, [v, w], config.dt)
        trajectory = np.vstack((trajectory, x))  # 垂直堆叠，vertical
        time += config.dt
    return trajectory


def calc_to_goal_cost(trajectory, goal, config):
    """
    计算轨迹到目标点的代价
    :param trajectory:轨迹搜索空间
    :param goal:
    :param config:
    :return: 轨迹到目标点欧式距离
    """
    # calc to goal cost. It is 2D norm.

    dx = goal[0] - trajectory[-1, 0]
    dy = goal[1] - trajectory[-1, 1]
    goal_dis = math.sqrt(dx**2 + dy**2)
    cost = config.to_goal_cost_gain * goal_dis

    return cost


def calc_obstacle_cost(traj, ob, config):
    """
    计算预测轨迹和障碍物的最小距离，dist(v,w)
    :param traj:
    :param ob:
    :param config:
    :return:
    """
    # calc obstacle cost inf: collision, 0:free

    min_r = float("inf")  # 距离初始化为无穷大

    for ii in range(0, len(traj[:, 1])):
        for i in range(len(ob[:, 0])):
            ox = ob[i, 0]
            oy = ob[i, 1]
            dx = traj[ii, 0] - ox
            dy = traj[ii, 1] - oy

            r = math.sqrt(dx**2 + dy**2)
            if r <= config.robot_radius:
                return float("Inf")  # collision

            if min_r >= r:
                min_r = r

    return 1.0 / min_r  # 越小越好 与障碍物距离越近，该值会越大；


def calc_final_input(x, u, vr, config, goal, ob):
    """
    计算采样空间的评价函数，选择最合适的那一个作为最终输入
    :param x:位置空间
    :param u:速度空间
    :param vr:速度空间交集
    :param config:
    :param goal:目标位置
    :param ob:障碍物
    :return:
    """
    x_init = x[:]
    min_cost = 10000.0
    min_u = u

    best_trajectory = np.array([x])
    for v in np.arange(vr[0], vr[1], config.v_reso):  
        for w in np.arange(vr[2], vr[3], config.yawrate_reso): 
            trajectory = calc_trajectory(x_init, v, w, config)

            # calc cost
            to_goal_cost = calc_to_goal_cost(
                trajectory, goal, config
            )  # trajectory与目标的欧式距离
            speed_cost = config.speed_cost_gain * (
                config.max_speed - trajectory[-1, 3]
            )  # v越大，该项越小，则效果越好；
            ob_cost = calc_obstacle_cost(
                trajectory, ob, config
            )  # 返回的是距离的倒数，该值约小，距离越大，越好
            final_cost = (
                to_goal_cost + speed_cost + 0 * ob_cost
            )  # 总计代价，越小轨迹约好

            # search minimum trajectory
            if min_cost >= final_cost:
                min_cost = final_cost
                min_u = [v, w]  
                best_trajectory = trajectory  

    return min_u, best_trajectory


def dwa_control(x, u, config, goal, ob):
    """
    调用前面的几个函数，生成最合适的速度空间和轨迹搜索空间
    :param x: 起始位置
    :param u:初始速度空间（v,w）
    :param config:
    :param goal:目标位置
    :param ob:障碍物
    :return:
    """
    # Dynamic Window control

    vr = calc_dynamic_window(x, config)  # 计算动态窗口，即机器人的电机实际物理参数范围

    u, trajectory = calc_final_input(
        x, u, vr, config, goal, ob
    )  # 选出最优轨迹搜索空间和速度空间

    return u, trajectory


def plot_arrow(x, y, yaw, length=0.5, width=0.1):
    plt.arrow(
        x,
        y,
        length * math.cos(yaw),
        length * math.sin(yaw),
        head_length=1.5 * width,
        head_width=width,
    )
    plt.plot(x, y)


def main():
    x = np.array([0.0, 0.0, math.pi / 2.0, 0.2, 0.0])

    goal = np.array([10, 10])

    # matrix二维矩阵 障碍物
    ob = np.matrix(
        [
            [-1, -1],
            [0, 2],
            [4.0, 2.0],
            [5.0, 4.0],
            [5.0, 5.0],
            [5.0, 6.0],
            [5.0, 9.0],
            [8.0, 9.0],
            [7.0, 9.0],
            [12.0, 12.0],
        ]
    )
    # ob = np.matrix([[0, 2]])
    # ob = np.matrix([[2, 2]])

    u = np.array([0.2, 0.0])  # 初始速度空间(v,w)
    config = Config()  # 初始化物理运动参数
    trajectory = np.array(x)  # 初始轨迹空间

    for i in range(1000):
        u, best_trajectory = dwa_control(
            x, u, config, goal, ob
        )  # 获取最优速度空间和轨迹参数
        x = motion(x, u, config.dt)  # 机器人运动方程 x为机器人状态空间

        trajectory = np.vstack((trajectory, x))  # store state history

        if show_animation:
            draw_dynamic_search(best_trajectory, x, goal, ob)

        # check goal
        if (
            math.sqrt((x[0] - goal[0]) ** 2 + (x[1] - goal[1]) ** 2)
            <= config.robot_radius
        ):
            print("Goal!!")
            break

    print("Done")

    draw_path(trajectory, goal, ob, x)


def draw_dynamic_search(best_trajectory, x, goal, ob):
    """
    画出动态搜索过程图
    :return:
    """
    plt.cla()  # 清除上次绘制图像
    plt.plot(best_trajectory[:, 0], best_trajectory[:, 1], "-g")
    plt.plot(x[0], x[1], "xr")
    plt.plot(0, 0, "og")
    plt.plot(goal[0], goal[1], "ro")
    plt.plot(ob[:, 0], ob[:, 1], "bs")
    plot_arrow(x[0], x[1], x[2])
    plt.axis("equal")
    plt.grid(True)
    plt.pause(0.0001)


def draw_path(trajectory, goal, ob, x):
    """
    画图函数
    :return:
    """
    plt.cla()  # 清除上次绘制图像

    plt.plot(x[0], x[1], "xr")
    plt.plot(0, 0, "og")
    plt.plot(goal[0], goal[1], "ro")
    plt.plot(ob[:, 0], ob[:, 1], "bs")
    plot_arrow(x[0], x[1], x[2])
    plt.axis("equal")
    plt.grid(True)
    plt.plot(trajectory[:, 0], trajectory[:, 1], "r")
    plt.show()


if __name__ == "__main__":
    main()
\end{minted}
